\subsection{INTEGRAL SUBGROUPS}

DEFINITION 1. Let G be a Lie group. An integral subgroup of G is a subgroup H with a connected Lie group structure such that the canonical injection of H into G is an immersion.

A one-parameter subgroup of \(G\) is a 1-dimensional integral subgroup of \(G\) .

Let H be an integral subgroup of G and i the canonical injection of H into G. Then \(L(i)\) defines an isomorphism of \(L(H)\) onto a Lie subalgebra of \(L(G)\) admitting a topological supplement. \(L(H)\) is identified with its image under \(L(i)\) .

Examples. (1) A connected Lie subgroup of \(\mathbf{G}\) is an integral subgroup of \(\mathbf{G}\) .

\begin{enumerate}
\def\labelenumi{(\arabic{enumi})}
\setcounter{enumi}{1}
\item
  Suppose that \(\mathbf{G}\) is finite-dimensional. Let \(\mathbf{H}\) be a subgroup of \(\mathbf{G}\) ; we give it the structure induced by the Lie group structure on \(\mathbf{G}\) (§ 4, no. 5, Definition 3). Then its identity component \(\mathrm{H}_0\) is an integral subgroup of \(\mathbf{G}\) and the subalgebra tangent to \(\mathbf{H}\) at \(e\) is \(\mathrm{L}(\mathrm{H}_0)\) (§ 4, no. 5, Proposition 9 (ii)).
\item
  Let \(\mathbf{G}\) be a complex Lie group, \(\mathbf{H}\) an integral subgroup of \(\mathbf{G}\) and \(\mathbf{G}_1\) (resp. \(\mathbf{H}_1\) ) the underlying real Lie group of \(\mathbf{G}\) (resp. \(\mathbf{H}\) ). Then \(\mathbf{H}_1\) is an integral subgroup of \(\mathbf{G}_1\) and \(\mathbf{L}(\mathbf{H}_1)\) is the underlying real Lie algebra of \(\mathbf{L}(\mathbf{H})\) .
\end{enumerate}

THEOREM 2. Let G be a Lie group.

\begin{enumerate}
\def\labelenumi{(\roman{enumi})}
\item
  The mapping \(\mathbf{H} \mapsto \mathbf{L}(\mathbf{H})\) is a bijection of the set of integral subgroups of \(\mathbf{G}\) onto the set of Lie subalgebras of \(\mathbf{L}(\mathbf{G})\) admitting a topological supplement.
\item
  Let \(\mathbf{H}\) be an integral subgroup of \(\mathbf{G}\) . Every connected Lie subgroup germ of \(\mathbf{G}\) with Lie algebra \(\mathbf{L}(\mathbf{H})\) is an open submanifold of \(\mathbf{H}\) which generates \(\mathbf{H}\) .
\end{enumerate}

\begin{enumerate}
\def\labelenumi{(\alph{enumi})}
\item
  Let \(\mathfrak{b}\) be a Lie subalgebra of \(\mathrm{L}(\mathrm{G})\) admitting a topological supplement. Let \(\mathrm{H}_{1}\) be a Lie subgroup germ of \(\mathrm{G}\) such that \(\mathrm{L}(\mathrm{H}_{1}) = \mathfrak{b}\) (§ 4, Theorem 3). \(\mathrm{H}_{1}\) can be chosen so that it is connected. Let \(\mathrm{H}\) be the subgroup of \(\mathrm{G}\) generated by \(\mathrm{H}_{1}\) . There exists (§ 1, Corollary to Proposition 22) a Lie group structure on \(\mathrm{H}\) such that \(\mathrm{H}_{1}\) is an open submanifold of \(\mathrm{H}\) and the canonical injection of \(\mathrm{H}\) into \(\mathrm{G}\) is an immersion. As \(\mathrm{H}_{1}\) is connected, \(\mathrm{H}\) is connected and is hence an integral subgroup of \(\mathrm{G}\) . Then \(\mathrm{L}(\mathrm{H}) = \mathrm{L}(\mathrm{H}_{1}) = \mathfrak{b}\) . This proves that the mapping considered in (i) is surjective.
\item
  Let H be an integral subgroup of G and \(N_{1}\) a connected Lie subgroup germ of G with Lie algebra L(H). As the canonical injection of H into G is an immersion, there exists an open subgroup germ \(H_{1}\) of H which is at the same time a submanifold of G and hence a Lie subgroup germ of G with Lie algebra L(H). On the other hand, let N be the subgroup of G generated by \(N_{1}\) ; by part (a) of the proof, it has the structure of an integral subgroup of G such that \(N_{1}\) is an open submanifold of N. By §4, Theorem 3, \(H_{1} \cap N_{1}\) is open in \(H_{1}\) and \(N_{1}\) . Hence the subgroup of G generated by \(H_{1} \cap N_{1}\) is equal on the one hand to H and on the other to N. Therefore the Lie groups H and N are equal. This proves (ii) and also proves that the mapping considered in (i) is injective.
\end{enumerate}

Remark 1. Let H be an integral subgroup of G. Let Y be the left foliation of G associated with L(H). If \(g \in G\) , let gH be given the manifold structure derived from that on H by \(\gamma(g)\) . By Differentiable and Analytic Manifolds, R, 9.3.2, the canonical injection of gH into Y is a morphism. This morphism is étale. Hence the maximal connected leaves of Y are the left cosets modulo H.

PROPOSITION 1. Let G and M be Lie groups, H an integral subgroup of G and \(\phi\) a morphism of M into G such that \(\mathrm{L}(\phi)(\mathrm{L}(\mathrm{M})) \subset \mathrm{L}(\mathrm{H})\) . Suppose that M is connected. Then \(\phi(\mathrm{M}) \subset \mathrm{H}\) and \(\phi\) , considered as a mapping of M into H, is a Lie group morphism.

In the notation of Remark 1, \(\phi\) is a morphism of M into Y (Differentiable and Analytic Manifolds, R, 9.3.2) and hence \(\phi(M) \subset H\) since M is connected.

COROLLARY 1. Let G and H be Lie groups, \(\phi\) a Lie group morphism of G into H, N the kernel of \(\phi\) and \(h = L(\phi)\) . Suppose that G is connected and that H is finite-dimensional.

\begin{enumerate}
\def\labelenumi{(\roman{enumi})}
\item
  N is a Lie subgroup of G and L(N) = Ker h.
\item
  Let \(\mathbf{H}'\) be the integral subgroup of \(\mathbf{H}\) with Lie algebra \(\operatorname{Im} h\) . Then \(\phi(\mathbf{G}) = \mathbf{H}'\) .
\item
  The mapping of G/N into H' derived from \(\phi\) when passing to the quotient is a Lie group isomorphism.
\end{enumerate}

\((i)\) has already been proved (§ 3, no. 8, Proposition 28).

Let \(\psi\) be the Lie group morphism of G/N into H derived from \(\phi\) when passing to the quotient; it is an immersion (§3, no. 8, Proposition 28). By Proposition 1, \(\psi\) is a Lie group morphism of G/N into H'. This morphism is étale and hence \(\psi\) (G/N) = H' since H' is connected; this proves (ii). Then \(\psi\colon G/N \to H'\) is bijective and is a Lie group isomorphism, which proves (iii).

COROLLARY 2. Let \(\mathbf{G}\) be a Lie group and \(\mathrm{H}_1\) and \(\mathrm{H}_2\) integral subgroups of \(\mathbf{G}\) . If \(\mathrm{L}(\mathrm{H}_2) \subset \mathrm{L}(\mathrm{H}_1)\) , then \(\mathrm{H}_2\) is an integral subgroup of \(\mathrm{H}_1\) .

Let \(i_1: \mathrm{H}_1 \to \mathrm{G}\) , \(i_2: \mathrm{H}_2 \to \mathrm{G}\) be the canonical injections. Then

\[
\mathrm{L} (i _ {2}) (\mathrm{L} (\mathrm{H} _ {2})) = \mathrm{L} (\mathrm{H} _ {2}) \subset \mathrm{L} (\mathrm{H} _ {1}).
\]

By Proposition 1, \(i_{2}\) is an analytic mapping of \(H_{2}\) into \(H_{1}\) and even an immersion of \(H_{2}\) into \(H_{1}\) since \(\mathrm{L}(i_{2})\) is an isomorphism of \(\mathrm{L}(\mathrm{H}_{2})\) onto a subalgebra of \(\mathrm{L}(\mathrm{H}_{1})\) admitting a topological supplement.

COROLLARY 3. Let \(\mathbf{G}\) be a finite-dimensional Lie group and \((\mathrm{H}_i)_{i\in I}\) a family of Lie subgroups of \(\mathbf{G}\) . Then \(\mathrm{H} = \bigcap_{i\in I}\mathrm{H}_i\) is a Lie subgroup of \(\mathbf{G}\) and

\[
\mathrm{L} (\mathrm{H}) = \bigcap_ {i \in \mathrm{I}} \mathrm{L} (\mathrm{H} _ {i}).
\]

There exists a finite subset J of I such that \(\bigcap_{i\in J}L(H_i)\) is equal to the intersection M of all the \(\mathrm{L}(\mathrm{H}_i)\) . We know that \(\mathrm{H}^* = \bigcap_{i\in J}\mathrm{H}_i\) is a Lie subgroup such that \(\mathrm{L}(\mathrm{H}^*) = \mathrm{M}\) (§ 3, no. 8, Corollary 2 to Proposition 29). Let \(\mathrm{H}_0\) be the identity component of \(\mathrm{H}^*\) . It is a Lie subgroup of G and \(\mathrm{L}(\mathrm{H}_0) = \mathrm{M}\) . By Corollary 2, \(\mathrm{H}_0\subset \mathrm{H}_i\) for all \(i\) and hence \(\mathrm{H}_0\subset \mathrm{H}\subset \mathrm{H}^*\) , whence the corollary.

COROLLARY 4. Let G be a finite-dimensional connected Lie group. The following conditions are equivalent:

\begin{enumerate}
\def\labelenumi{(\roman{enumi})}
\item
  G is unimodular (Integration, Chapter VII, § 1, no. 3, Definition 3);
\item
  det Ad g = 1 for all g ∈ G;
\item
  \(\operatorname{Tr} \operatorname{ad} a = 0\) for all \(a \in \mathrm{L}(\mathrm{G})\) .
\end{enumerate}

The mapping \(g \mapsto \operatorname{det} \operatorname{Ad} g\) is a morphism \(\phi\) of \(G\) into \(K^*\) . By § 3, Propositions 35 (no. 10) and 44 (no. 12), \(L(\phi)a = Tr\operatorname{ad} a\) for all \(a \in L(G)\) . Clearly \(\operatorname{Im} L(\phi) = \{0\}\) or \(K\) . In the first (resp. second) case, \(\operatorname{Im} \phi = \{1\}\) (resp. \(\operatorname{Im} \phi = K^*\) ) by Corollary 1 and hence \(G\) is unimodular (resp. not unimodular) by § 3, no. 16, Corollary to Proposition 55.

PROPOSITION 2. Let \(\mathbf{G}\) be a finite-dimensional Lie group and \(\mathbf{H}\) an integral subgroup of \(\mathbf{G}\) . The following conditions are equivalent:

\begin{enumerate}
\def\labelenumi{(\roman{enumi})}
\item
  H is closed;
\item
  the topology on H is induced by that on G;
\item
  H is a Lie subgroup of G.
\end{enumerate}

\((i) \Rightarrow (iii)\): this follows from § 1, Propositions 2 (iv) (no. 1) and 14 (iii) (no. 7).

\((iii) \Rightarrow (ii)\): obvious.

\((ii) \Rightarrow (i)\): if the topology on H is induced by that on G, H is closed because H is complete (§ 1, no. 1, Proposition 1).

PROPOSITION 3. Let G be a Lie group, H an integral subgroup of G, M a non-empty connected analytic manifold, f a mapping of M into G and \(r \in N_{K}\) . Consider the following conditions:

\begin{enumerate}
\def\labelenumi{(\roman{enumi})}
\item
  \(f\) is of class \(\mathbf{C}^r\) and \(f(\mathbf{M})\subset \mathbf{H}\) ;
\item
  \(f(\mathbf{M})\subset \mathbf{H}\) and \(f\) , considered as a mapping of \(\mathbf{M}\) into \(\mathbf{H}\) , is of class \(\mathbf{C}^r\)
\item
  \(f\) is of class \(\mathbf{C}^r\) , \(f(\mathbf{M})\) meets \(\mathbf{H}\) and the image of \(\mathrm{T}_m(\mathbf{M})\) is contained in \(f(m)\) . \(\mathbf{L}(\mathbf{H})\) for all \(m \in \mathbf{M}\) .
\end{enumerate}

Then (ii) \(\Leftrightarrow\) (iii) \(\Rightarrow\) (i). If the topology on H admits a countable base, the three conditions are equivalent.

\((ii) \Rightarrow (i)\) and \((ii) \Rightarrow (iii)\): obvious.

\((iii) \Rightarrow (ii)\): suppose that condition (iii) holds. By Differentiable and Analytic Manifolds, R, 9.2.8, \(f\) is a morphism of class \(\mathbf{C}^r\) of M into the left foliation associated with \(\mathbf{L}(\mathbf{H})\) . As M is connected, \(f(\mathbf{M}) \subset \mathbf{H}\) .

If the topology on H admits a countable base, condition (i) implies that f is a mapping of class \(C^{r}\) of M into H (Differentiable and Analytic Manifolds, R, 9.2.8); hence (i) \(\Rightarrow\) (ii).

COROLLARY 1. Let G be a finite-dimensional Lie group and H an integral subgroup of G. Then the Lie subalgebra tangent to H at e (§ 4, no. 5, Definitions 2 and 3) is L(H) and the Lie group structure on H is the structure induced by that on G.

As H is connected and finite-dimensional, its topology admits a countable base.

COROLLARY 2. Let G be a Lie group and \(H_{1}\) and \(H_{2}\) integral subgroups of G. Suppose that the topology on \(H_{1}\) admits a countable base. Then

\[
\mathrm{H} _ {2} \subset \mathrm{H} _ {1} \Leftrightarrow \mathrm{L} (\mathrm{H} _ {2}) \subset \mathrm{L} (\mathrm{H} _ {1})
\]

and, if these conditions hold, \(\mathbf{H}_2\) is an integral subgroup of \(\mathbf{H}_1\) .

The last assertion and the implication \(\mathrm{L}(\mathrm{H}_{2})\subset\mathrm{L}(\mathrm{H}_{1})\Rightarrow\mathrm{H}_{2}\subset\mathrm{H}_{1}\) follow from Corollary 2 to Proposition 1. The converse implication follows from Proposition 3.

COROLLARY 3. Let G be a Lie group and \(H_{1}\) and \(H_{2}\) integral subgroups of G whose topology admits a countable base. If \(H_{1}\) and \(H_{2}\) have the same underlying set, the Lie group structures on \(H_{1}\) and \(H_{2}\) are the same.

This follows from Corollary 2.

Remark 2. Let G be a finite-dimensional Lie group. Let H be a subgroup of G. We shall say, by an abuse of language, that H is an integral subgroup of G if there exists a Lie group structure S on H such that H, together with S, is an integral subgroup of G. By Corollary 3 to Proposition 3, if S exists, S is unique.

Remark 3. Let V be a manifold of class \(\mathbf{C}^r\) . Let M be a subset of V and \(x\) and \(y\) elements of M. Consider the following property:

\(P_{M,x,y}\) : there exist I, \(x_{0}, x_{1}, \ldots, x_{n}, f_{1}, \ldots, f_{n}\) such that: (a) I is a connected open subset of K; (b) \(x_{0}, \ldots, x_{n}\) are in M, \(x_{0} = x, x_{n} = y\) ; (c) for \(1 \leqslant i \leqslant n\) , \(f_{i}\) is a mapping of class \(C^{r}\) of I into V which takes the values \(x_{i-1}\) and \(x_{i}\) and \(f_{i}(I) \subset M\) .

We shall say that M is a \(C^{r}\) -connected subset of V if, for all elements x, y of M, the property \(P_{M,x,y}\) holds.

PROPOSITION 4. Let G be a finite-dimensional Lie group and H a subgroup of G. Let \(r \in \mathbf{N}_{\mathbb{K}}\) . The following conditions are equivalent:

\begin{enumerate}
\def\labelenumi{(\roman{enumi})}
\item
  H is an integral subgroup of G;
\item
  with the Lie group structure induced by that on G, H is connected;
\item
  H is \(\mathbf{C}^r\) -connected.
\end{enumerate}

\((ii) \Rightarrow (i)\): obvious.

\((i) \Rightarrow (iii)\): suppose that H has a Lie group structure such that H is an integral subgroup of G. Using the notation of Remark 3, the set of \(y \in H\) such that property \(\mathrm{P}_{\mathrm{H},e,y}\) holds is an open subgroup of H. As H is connected, this subgroup is equal to H and hence condition (iii) is satisfied.

\((iii) \Rightarrow (ii)\): suppose that condition (iii) is satisfied and let H be given the structure induced by the Lie group structure on G. Let \(\mathfrak{h}\) be the subalgebra tangent to H at \(e\) . The identity component \(\mathrm{H}_0\) of H is an integral subgroup of G such that \(\mathrm{L}(\mathrm{H}_0) = \mathfrak{h}\) . We show that \(\mathrm{H} = \mathrm{H}_0\) . It suffices to prove the following: let I be a connected open subset of K, \(f\) a mapping of class \(\mathbf{C}^r\) of I into G such that \(f(\mathrm{I}) \subset \mathrm{H}\) and \(\lambda\) and \(\mu\) two points of I; if \(f(\lambda) \in \mathrm{H}_0\) , then \(f(\mu) \in \mathrm{H}_0\) . But, for all \(\nu \in \mathrm{I}\) , \((\mathrm{T}_{\nu}f)(\mathrm{K}) \subset f(\nu)\mathfrak{h}\) by definition of \(\mathfrak{h}\) , so that our assertion follows from Proposition 3.

Remark 4. If \(\mathbf{K} = \mathbf{R}\) , the integral subgroups of \(\mathbf{G}\) can also be characterized as the subgroups which, with the topology induced by that on \(\mathbf{G}\) , are arcwise connected (§ 8, Exercise 4). However, subgroups may be connected but not integral (Commutative Algebra, Chapter VI, § 9, Exercise 2).

COROLLARY. Let G be a finite-dimensional Lie group and \(H_{1}\) and \(H_{2}\) two integral subgroups of G. The subgroup of G generated by \(H_{1}\) and \(H_{2}\) and the subgroup ( \(H_{1}, H_{2}\) ) of G are integral subgroups of G.

The subgroup (G, G) of G is not always closed (§ 9, Exercise 6).

Recall (§ 3, no. 11, Corollary 5 to Proposition 41) that if a is a finite-dimensional algebra, \(\operatorname{Aut}(a)\) is a Lie subgroup of \(\mathbf{GL}(a)\) and that \(\mathrm{L}(\mathrm{Aut}(a))\) is the Lie algebra of derivations of \(a\) .

DEFINITION 2. Let \(\mathfrak{a}\) be a finite-dimensional Lie algebra. Let \(\operatorname{Ad}(\mathfrak{a})\) or \(\operatorname{Int}(\mathfrak{a})\) denote the integral subgroup of \(\operatorname{Aut}(\mathfrak{a})\) with Lie algebra \(\operatorname{ad}(\mathfrak{a})\) . The elements of this group are called inner automorphisms of \(\mathfrak{a}\) .

By transport of structure, \(\operatorname{ad}(a)\) is invariant under \(\operatorname{Aut}(a)\) and hence \(\operatorname{Int}(a)\) is normal in \(\operatorname{Aut}(a)\) . By § 4, no. 4, Corollary 1 to Proposition 8 and the fact that \(\operatorname{Int}(a)\) is connected, the elements of \(\operatorname{Int}(a)\) are the finite products of automorphisms of the form \(\exp ad x\) where \(x \in a\) . In general, \(\operatorname{Int}(a)\) is not a Lie subgroup of \(\operatorname{Aut}(a)\) (Exercise 14).

